% !TEX root = ../main.tex
\begin{ouverture}[spectre]{19}{Couleurs et vision}{Partie 4 · Ondes et signaux}
{Lumière blanche\enspace·\enspace Synthèses additive et soustractive\enspace·\enspace Couleur des objets\enspace·\enspace Filtres\enspace·\enspace Trichromie}
\objectif{Choisir le modèle de la synthèse additive ou celui de la synthèse soustractive selon la situation.}
\objectif{Interpréter la couleur perçue d'un objet à partir de la lumière incidente et des phénomènes d'absorption, de diffusion et de transmission.}
\objectif{Prévoir le résultat de la superposition de lumières colorées et l'effet de filtres colorés.}
\objectif{Relier la vision des couleurs aux trois types de cônes de la rétine.}
\prerequis{Lumière blanche, spectre, dispersion par un prisme, longueur d'onde (classe de seconde).}
\prerequis{Absorbance et couleur d'une solution (chapitre 2).}
\end{ouverture}
\essentiel
\essSection{1}{Lumière blanche et couleurs complémentaires}
\begin{definition}[Lumière blanche et couleurs]
La lumière blanche contient toutes les radiations visibles (de \qty{400}{nm} à \qty{800}{nm} environ). Elle peut être reproduite en superposant
trois lumières colorées dites \textbf{primaires} : rouge, verte et bleue. Deux lumières colorées sont \textbf{complémentaires} si leur
superposition donne du blanc : rouge et cyan, vert et magenta, bleu et jaune.
\end{definition}
\schema{19-cercle}
\essSection{2}{Synthèse additive et synthèse soustractive}
\begin{propriete}[Deux modèles]
\textbf{Synthèse additive} : on superpose des \textbf{lumières} ; les couleurs s'ajoutent. Rouge + vert = jaune, vert + bleu = cyan, rouge + bleu =
magenta, rouge + vert + bleu = blanc. Exemples : écrans, projecteurs de spectacle.\par
\textbf{Synthèse soustractive} : on superpose des \textbf{filtres} ou des encres, qui retirent chacun une partie de la lumière blanche. Un filtre jaune
retire le bleu, un cyan le rouge, un magenta le vert ; les trois ensemble ne laissent rien passer (noir). Exemples : imprimerie,
peinture, photographie argentique.
\end{propriete}
\schema{19-synthese}
\essSection{3}{Couleur des objets}
\begin{loi}[Absorption, diffusion, transmission]
Un objet opaque \textbf{absorbe} une partie de la lumière qu'il reçoit et \textbf{diffuse} le reste dans toutes les directions : sa couleur perçue est
celle de la lumière diffusée. Un filtre ou une solution colorée absorbe une partie de la lumière et \textbf{transmet} le reste. La couleur
perçue dépend donc aussi de la lumière incidente : une tomate, qui diffuse le rouge, paraît noire en lumière verte.
\end{loi}
\schema{19-objet}
\begin{python}
blanche = {'rouge', 'vert', 'bleu'}       # lumiere blanche
filtre_jaune = {'rouge', 'vert'}          # transmet rouge et vert
filtre_cyan = {'vert', 'bleu'}            # transmet vert et bleu
print(sorted(blanche & filtre_jaune))
print(sorted(blanche & filtre_jaune & filtre_cyan))
\end{python}
\begin{console}
['rouge', 'vert']
['vert']
\end{console}
\noindent\capacite{Capacité numérique} Une lumière est représentée par l'ensemble de ses composantes ; un filtre ne garde que les composantes qu'il
transmet (intersection \code{\&}). Deux filtres jaune et cyan ne laissent passer que le vert.
\essSection{4}{Vision des couleurs}
\begin{propriete}[Trichromie]
La rétine contient trois types de \textbf{cônes}, sensibles surtout au bleu, au vert ou au rouge. Le cerveau reconstitue la couleur à partir des
signaux de ces trois capteurs : c'est la \textbf{trichromie}. Un écran trompe l'œil en n'émettant que trois lumières, rouge, verte et bleue,
d'intensités réglables. L'absence d'un type de cônes entraîne un daltonisme.
\end{propriete}
\schema{19-cones}
\begin{attention}
Une peinture jaune ne produit pas de lumière jaune : elle absorbe le bleu et diffuse le rouge et le vert. Mélanger des peintures relève
de la synthèse soustractive, superposer des lumières de la synthèse additive : bleu + jaune donne du blanc en lumière, mais du vert en
peinture (cyan et jaune, en pratique).
\end{attention}
\begin{experience}[Filtres, projecteurs et objets colorés]
Avec trois projecteurs munis de filtres rouge, vert et bleu, on éclaire un écran blanc puis des objets colorés. On observe les zones de
superposition, puis la couleur des objets sous différents éclairages, et l'effet de filtres placés devant une lampe blanche.
\end{experience}
\begin{rediger}
Préciser le modèle utilisé (lumières qui s'ajoutent ou filtres qui retranchent), décomposer chaque lumière en rouge, vert et bleu,
puis indiquer les composantes absorbées, diffusées ou transmises avant de conclure sur la couleur perçue.
\end{rediger}
\exercices

\rubrique{Découverte}
\begin{blocLivre}{ex-19-1}
\begin{exercice}[1]{Trois projecteurs}
\schema{19-ex01}
\begin{questions}
\item Quel modèle de synthèse faut-il utiliser ? Justifier.
\item Indiquer la couleur des zones 1, 2, 3 et 4.
\item On éteint le projecteur bleu. Quelle est alors la couleur de la zone 4 ?
\item Avec deux projecteurs seulement, comment obtenir du blanc ?
\end{questions}
\end{exercice}
\end{blocLivre}
\begin{blocLivre}{ex-19-2}
\begin{exercice}[1]{Trois filtres superposés}
\schema{19-ex02}
\begin{questions}
\item Quel modèle de synthèse faut-il utiliser ? Justifier.
\item Quelle couleur chaque filtre absorbe-t-il ?
\item Indiquer la couleur des zones 1, 2, 3 et 4.
\item Citer une technique qui utilise ce type de synthèse.
\end{questions}
\end{exercice}
\end{blocLivre}
\begin{blocLivre}{ex-19-3}
\begin{exercice}[1]{La couleur d'une tomate}
Éclairée en lumière blanche, une tomate paraît rouge.
\begin{questions}
\item Quelles lumières la tomate absorbe-t-elle ? Laquelle diffuse-t-elle ?
\item Quelle est sa couleur en lumière verte ? En lumière jaune ? En lumière cyan ?
\item Un citron paraît jaune en lumière blanche. Quelles lumières diffuse-t-il ?
\item Quelle est la couleur du citron en lumière rouge ? En lumière bleue ?
\end{questions}
\end{exercice}
\end{blocLivre}
\begin{blocLivre}{ex-19-4}
\begin{exercice}[1]{Des filtres successifs}
\schema{19-ex04}
\begin{questions}
\item Quelles lumières chaque filtre transmet-il ?
\item Quelle lumière sort du premier montage ? Du second ?
\item On remplace les deux filtres par un filtre magenta suivi d'un filtre jaune. Quelle lumière sort ?
\item Une lumière magenta éclaire un filtre cyan. Quelle lumière traverse ?
\end{questions}
\end{exercice}
\end{blocLivre}
\begin{blocLivre}{ex-19-5}
\begin{exercice}[1]{Couleurs complémentaires}
\schema{19-cercle}
\begin{questions}
\item Citer les trois couples de couleurs complémentaires.
\item Que donne la superposition de lumières rouge et cyan ?
\item Une solution absorbe surtout le vert. De quelle couleur paraît-elle ? Et une solution qui absorbe surtout le bleu ?
\item Pourquoi parle-t-on de couleur « complémentaire » ?
\end{questions}
\end{exercice}
\end{blocLivre}
\rubrique{Entraînement}
\begin{blocLivre}{ex-19-6}
\begin{exercice}[2]{Des objets sous un projecteur jaune, en Python}
\begin{python}
couleurs = {'rouge': {'R'}, 'vert': {'V'}, 'bleu': {'B'},
            'jaune': {'R', 'V'}, 'cyan': {'V', 'B'}, 'magenta': {'R', 'B'},
            'blanc': {'R', 'V', 'B'}}
lumiere = couleurs['jaune']
for objet in ['blanc', 'rouge', 'cyan', 'bleu']:
    vu = lumiere & couleurs[objet]
    nom_vu = [n for n in couleurs if couleurs[n] == vu]
    print(objet, '->', nom_vu[0] if nom_vu else 'noir')
\end{python}
\begin{console}
blanc -> jaune
rouge -> rouge
cyan -> vert
bleu -> noir
\end{console}
\begin{questions}
\item \capacite{Capacité numérique}Que représente l'ensemble associé à un objet ? Que calcule \code{lumiere \& couleurs[objet]} ?
\item Justifier chacune des quatre lignes affichées.
\item Quelle est la couleur d'un objet magenta sous ce projecteur ?
\item Comment modifier le programme pour un projecteur cyan ? Prévoir le résultat pour les quatre objets.
\end{questions}
\end{exercice}
\end{blocLivre}
\begin{blocLivre}{ex-19-7}
\begin{exercice}[2]{Les couleurs d'un écran}
Chaque pixel d'un écran contient trois sous-pixels rouge, vert et bleu. L'intensité de chacun est codée par un nombre entier de 0 (éteint)
à 255 (intensité maximale) : une couleur est notée $(R ; V ; B)$.
\begin{questions}
\item Quel type de synthèse un écran utilise-t-il ?
\item Indiquer la couleur obtenue pour $(255 ; 255 ; 0)$, $(255 ; 0 ; 255)$, $(0 ; 255 ; 255)$, $(255 ; 255 ; 255)$ et $(0 ; 0 ; 0)$.
\item Que donne $(128 ; 128 ; 128)$ ?
\item Combien de couleurs différentes l'écran peut-il afficher ?
\end{questions}
\end{exercice}
\end{blocLivre}
\begin{blocLivre}{ex-19-8}
\begin{exercice}[2]{Les éclairages d'un spectacle}
Sur scène, deux projecteurs, l'un rouge et l'autre bleu, éclairent un danseur vêtu de blanc et une danseuse en robe rouge.
\begin{questions}
\item De quelle couleur paraît le costume blanc ? Justifier.
\item De quelle couleur paraît la robe rouge ?
\item On éteint les deux projecteurs et l'on allume un projecteur vert. Décrire l'aspect des deux costumes.
\item Quel éclairage faut-il pour qu'un costume jaune paraisse rouge ?
\end{questions}
\end{exercice}
\end{blocLivre}
\begin{blocLivre}{ex-19-9}
\begin{exercice}[2]{Spectres d'absorption}
\schema{19-ex09}
\begin{questions}
\item Quelles radiations la solution A absorbe-t-elle le plus ? Et la solution B ?
\item En déduire la couleur de chaque solution éclairée en lumière blanche.
\item Quel modèle de synthèse utilise-t-on pour interpréter ces couleurs ?
\item On éclaire la solution A avec une lumière rouge. Que voit-on ?
\end{questions}
\end{exercice}
\end{blocLivre}
\begin{blocLivre}{ex-19-10}
\begin{exercice}[2]{Les cônes et le daltonisme}
\schema{19-cones}
\begin{questions}
\item Combien de types de cônes la rétine contient-elle ? À quelles couleurs chacun est-il le plus sensible ?
\item Une lumière jaune de \qty{580}{nm} excite-t-elle un seul type de cônes ? Comment le cerveau l'identifie-t-il ?
\item Pourquoi un écran peut-il donner l'illusion du jaune avec seulement du rouge et du vert ?
\item Une personne n'a pas de cônes L. Pourquoi confond-elle le rouge et le vert ?
\end{questions}
\end{exercice}
\end{blocLivre}
\begin{blocLivre}{ex-19-11}
\begin{exercice}[2]{Le capteur d'un appareil photo}
Chaque photosite du capteur d'un appareil photo ne mesure que l'intensité lumineuse. Pour obtenir des couleurs, on place devant lui une
mosaïque de filtres rouges, verts et bleus : sur un carré de quatre photosites, deux sont verts, un est rouge, un est bleu.
\begin{questions}
\item Quel est le rôle des filtres colorés ? Quel modèle de synthèse interviennent-ils ?
\item Pour un capteur de 12 millions de photosites, combien sont équipés d'un filtre vert ? D'un filtre rouge ?
\item Pourquoi place-t-on deux fois plus de filtres verts ? Utiliser le graphique des cônes.
\item Comment l'appareil reconstitue-t-il la couleur de chaque pixel de l'image ?
\end{questions}
\end{exercice}
\end{blocLivre}
\rubrique{Objectif DS}
\begin{blocLivre}{ex-19-12}
\begin{exercice}[3]{L'imprimante}
\begin{docu}{1}{Les cartouches}
Une imprimante utilise des encres cyan, magenta et jaune, déposées en couches transparentes sur une feuille blanche, ainsi qu'une
encre noire.
\end{docu}
\begin{questions}
\item Quel modèle de synthèse décrit l'impression ? Justifier.
\item Quelles encres faut-il superposer pour imprimer du bleu ? Du rouge ? Du vert ?
\item \initiative On regarde l'impression en lumière rouge. Comment apparaissent les zones cyan, jaune et rouge ?
\item En théorie, les trois encres superposées donnent du noir. Pourquoi ajoute-t-on une encre noire ?
\end{questions}
\end{exercice}
\end{blocLivre}
\begin{blocLivre}{ex-19-13}
\begin{exercice}[3]{Les lunettes anaglyphes}
\begin{docu}{1}{Voir en relief}
Une image « en relief » est imprimée sur papier blanc en deux exemplaires légèrement décalés : l'un rouge, l'autre cyan. On la regarde
avec des lunettes munies d'un filtre rouge devant l'œil gauche et d'un filtre cyan devant l'œil droit. Le cerveau combine les deux
images reçues et crée l'impression de relief.
\end{docu}
\begin{questions}
\item À travers le filtre rouge, de quelle couleur paraissent le fond blanc, le dessin rouge et le dessin cyan ?
\item Même question à travers le filtre cyan.
\item \initiative Expliquer pourquoi chaque œil ne voit qu'un seul des deux dessins.
\item Pourquoi utilise-t-on deux couleurs complémentaires ?
\end{questions}
\end{exercice}
\end{blocLivre}
\begin{blocLivre}{ex-19-14}
\begin{exercice}[3]{Des pixels invisibles}
\begin{docu}{1}{L'œil}
L'œil distingue deux points séparés d'un angle d'au moins une minute d'angle ($\dfrac{1}{60}$ de degré), soit \qty{2,9e-4}{rad} ; la distance entre les
points vaut alors environ cet angle multiplié par la distance d'observation.
\end{docu}
\begin{docu}{2}{Deux écrans}
Un smartphone affiche 400 pixels par pouce (\qty{1}{pouce}~$=\qty{25,4}{mm}$) ; une télévision de \qty{1,0}{m} de large affiche 1920 pixels sur sa largeur.
\end{docu}
\begin{questions}
\item Calculer la taille d'un pixel du smartphone.
\item Quel est le plus petit détail que l'œil distingue à \qty{30}{cm} ? Voit-on les pixels du smartphone ?
\item \initiative À partir de quelle distance les sous-pixels de la télévision ne sont-ils plus distingués ?
\item Expliquer pourquoi l'on voit des couleurs variées et non des points rouges, verts et bleus.
\end{questions}
\end{exercice}
\end{blocLivre}
\begin{blocLivre}{ex-19-15}
\begin{exercice}[3]{La couleur des feuilles}
\begin{docu}{1}{La chlorophylle}
La chlorophylle, pigment des feuilles, absorbe fortement le bleu et le rouge et très peu le vert. L'énergie de la lumière absorbée
permet la photosynthèse. En automne, la chlorophylle se dégrade ; il reste des pigments, les caroténoïdes, qui absorbent surtout le bleu.
\end{docu}
\begin{docu}{2}{Une serre éclairée}
Certaines serres sont éclairées par des LED rouges et bleues seulement.
\end{docu}
\begin{questions}
\item Expliquer la couleur verte des feuilles en lumière blanche.
\item Expliquer leur couleur en automne.
\item \initiative Pourquoi éclaire-t-on les serres en rouge et en bleu, et pas en vert ?
\item De quelle couleur paraissent les plantes dans cette serre ? Justifier.
\end{questions}
\end{exercice}
\end{blocLivre}
